\subsection{CT11 --- Localization}\label{ct:11}\label{additivity-and-the-localization-lemma}
\ctsignature{\(P_{5\to11}\to \ctcert{\ctnone};\ \ctint{P_{11\to1},P_{11\to7},P_{11\to14}};\ \ctrec{P_{11\to10}}.\)}

CT11 is a route-only localization machine.  Its entry contract records the
global deficit, admissible decomposition, admissibility class, and bounded
local type.  The executable framework adds four independent propositions:
decomposition certification, closure of admissibility under local moves,
summation, and existence of a localized deficit.  It also fixes the exact
validated entry interfaces of CT1, CT7, CT10, and CT14.

The core plan has five fields: \texttt{scope}, \texttt{decomposition},
\texttt{admissibility}, \texttt{localization}, and \texttt{routing}.  The
decomposition node and localization node each have one successor and therefore
produce certificates rather than choices.  The admissibility decision either
constructs an exact CT10 input or proves that the localized object remains in
the admissible class.  Only the latter state may enter localization.  The final
routing decision constructs an exact CT1 target-test input, CT7 exchange input,
or CT14 aggregate-mass input.  CT11 has no closure constructor.

This separation makes two common mistakes unrepresentable.  An uncertified
decomposition cannot reach the admissibility decision, and an admissibility gap
cannot be silently carried into the averaging lemma.  Each handoff structure
contains a downstream input and an alignment proof tying it to the CT11 entry.

\begin{itemize}
\tightlist
\item \textbf{S-Def} is discharged by the validated entry and \texttt{DecompositionState}.
\item \textbf{S-Comp} is discharged by \texttt{LocalizationState}, which stores both summation and localized-deficit proofs.
\item \textbf{S-Rout} and \textbf{S-Trig} are discharged by the four exact consumer payloads.
\end{itemize}

\begin{figure}[p]
\centering
\resizebox{0.98\textwidth}{!}{%
\begin{tikzpicture}[x=1cm,y=1cm,every node/.style={font=\scriptsize}]
\node[bcwide] (entry) at (0,0) {{\tiny\ttfamily CT11.entry}\\\textbf{Validated deficit and decomposition}};
\node[bcdec,bclemma] (scope) at (0,-2.2) {{\tiny\ttfamily CT11.decide.scope}\\\textbf{Finite local type?}};
\node[bcscopefail] (scopeout) at (-6.2,-2.2) {{\tiny\ttfamily CT11.terminal.scope}\\\textbf{Scope exit}};
\node[bcwide,bclemma] (decomp) at (0,-4.5) {{\tiny\ttfamily CT11.certify.decomposition}\\\textbf{Admissible decomposition}};
\node[bcdec,bclemma] (admit) at (0,-6.9) {{\tiny\ttfamily CT11.decide.\allowbreak admissibility}\\\textbf{Closed under local moves?}};
\node[bcroute,bcrouteneutral] (ct10) at (6.4,-6.9) {{\tiny\ttfamily CT11.terminal.ct10}\\\textbf{CT 10 exit}};
\node[bcwide,bclemma] (localize) at (0,-9.4) {{\tiny\ttfamily CT11.certify.localization}\\\textbf{Summation and local deficit}};
\node[bcroutechoice,bctheorem] (routing) at (0,-11.9) {{\tiny\ttfamily CT11.decide.routing}\\\textbf{Classify localized object}};
\node[bcroute,bcrouteneutral] (ct1) at (-6.2,-14.5) {{\tiny\ttfamily CT11.terminal.ct1}\\\textbf{CT 1 exit}};
\node[bcroute,bcrouteneutral] (ct7) at (0,-14.5) {{\tiny\ttfamily CT11.terminal.ct7}\\\textbf{CT 7 exit}};
\node[bcroute,bcrouteneutral] (ct14) at (6.2,-14.5) {{\tiny\ttfamily CT11.terminal.ct14}\\\textbf{CT 14 exit}};
\draw[bcarr] (entry)--(scope);
\draw[bcscopearr] (scope)--node[bclabel,above]{exit}(scopeout);
\draw[bcarr] (scope)--node[bclabel,right]{ready}(decomp);
\draw[bcarr] (decomp)--(admit);
\draw[bchandoff] (admit)--node[bclabel,above]{gap}(ct10);
\draw[bcarr] (admit)--node[bclabel,right]{closed}(localize);
\draw[bcarr] (localize)--(routing);
\draw[bchandoff] (routing)--node[bclabel,above,sloped]{target tests}(ct1);
\draw[bchandoff] (routing)--node[bclabel,right]{exchange}(ct7);
\draw[bchandoff] (routing)--node[bclabel,above,sloped]{mass}(ct14);
\end{tikzpicture}%
}
\caption[Closure-tactic 11: exact localization machine]{CT11 as an exact route-only machine.  Certification of the decomposition, the admissibility decision, and certification of localization are distinct interfaces.}
\label{fig:ct11-localization}
\end{figure}
\FloatBarrier
